\documentclass[../../../include/open-logic-section]{subfiles}

\begin{document}

\olfileid[pt]{sth}{replacement}{limofsize}
\olsection{Limitação de Tamanho}

Talvez a tentativa mais comum de oferecer uma justificação ``intrínseca'' da
Substituição passe pela seguinte noção:
\begin{enumerate}
	\item[] \limofsize. Quaisquer coisas formam um conjunto, desde que não haja
	demasiadas.
\end{enumerate}
Este princípio vindicará de imediato a Substituição. Afinal, qualquer
conjunto formado por Substituição não pode ser maior do que o conjunto a partir do qual
foi formado. Enunciado com precisão: suponha que forma um conjunto
$\funimage{\tau}{A} = \Setabs{\tau(x)}{x \in A}$ usando Substituição;
então $\cardle{\funimage{\tau}{A}}{A}$; logo, se os !!{element}os de $A$
não eram demasiado numerosos para formar um conjunto, as suas imagens não são demasiado numerosas
para formar $\funimage{\tau}{A}$. 

A dificuldade óbvia em invocar \limofsize{} para justificar a
Substituição é que \emph{ainda} não estabelecemos qualquer princípio
como \limofsize. Além disso, quando contamos a nossa história sobre a
concepção cumulativa-iterativa de conjunto em
\crefrange{sth:story::chap}{sth:z::chap}, nada \emph{sugeriu}
a direção de \limofsize. Isto é, de fato, precisamente a razão pela qual Boolos
escreveu num dado momento: ``Talvez se possa concluir que há pelo menos
dois pensamentos `por detrás' da teoria de conjuntos'' \citeyearpar[p.~19]{Boolos1989}. Por um lado, as ideias em torno da concepção cumulativa-iterativa
de conjunto destinam-se a vindicar~$\Z$. Por outro,
\limofsize{} visa vindicar a Substituição. 

Mas a questão não é apenas termos estado até aqui \emph{em silêncio}
sobre \limofsize. Antes, a questão é que \limofsize{} (tal como acabamos de o
formular) parece assentar mal com a noção cumulativa-iterativa
de conjunto. Afinal, não menciona nada sobre a ideia de conjuntos
formados em \emph{estágios}.

Isto não surpreende muito, dada a história destes
``dois pensamentos'' (O.s.\@, a concepção cumulativa-iterativa de conjunto,
e \limofsize). Estes ``dois pensamentos'' acabam por ser dois
projetos bastante diferentes para bloquear os paradoxos da teoria de conjuntos.
A noção cumulativa-iterativa de conjunto bloqueia o paradoxo de Russell
dizendo, em traços largos: \emph{nunca deveríamos ter esperado que existisse um conjunto de Russell,
porque não seria ``formado'' em nenhum estágio}. Por contraste,
\limofsize{} visa excluir o conjunto de Russell, dizendo, em traços largos:
\emph{nunca deveríamos ter esperado que existisse um conjunto de Russell,
porque seria demasiado grande}. 

Posto assim, sejamos diretos: considerada como resposta aos
paradoxos, \limofsize{} carece de \emph{muito} mais
justificação. Considere, por exemplo, esta versão do paradoxo de Russell: \emph{nenhum pug fareja exatamente os pugs que não se cheiram
a si próprios} (ver \olref[sth][story][rus]{sec}). Se perguntar ``por que não há tal pug?'', não é boa resposta dizer-lhe que tal pug teria de cheirar demasiados
pugs. Logo por que seria boa explicação intuitiva, da
não existência de um conjunto de Russell, que teria de ser ``demasiado grande''
para existir? 

Em suma, é perdoável estar um pouco perplexo quanto à motivação ``intuitiva''
para \limofsize. 

\end{document}
